\documentclass[10pt]{article}
\usepackage[margin=0.7in]{geometry}
\usepackage{amsmath}
\usepackage{booktabs}
\usepackage{array}
\usepackage{hyperref}
\usepackage{graphicx}
\usepackage{pdflscape}
\usepackage{float}

\title{AC-OPF Surrogates: GridFM, GridSFM, and LUMINA\\
\large Two data pipelines, two training objectives, six grids}
\author{PIGNN-Attn-LS-PPC experimental summary}
\date{8 August 2026}

\begin{document}
\maketitle

\section{Scope}
This report compares three power-system OPF surrogates across \textbf{two
independently constructed data pipelines} and, on OPFData, \textbf{two training
objectives}.

\paragraph{What the three models are --- they are not alike.} LUMINA and GridSFM
are \emph{released foundation models} with public checkpoints, evaluated here
zero-shot, fine-tuned and from scratch. \textbf{GridFM is not.} The ``GridFM'' in
this study is a local re-implementation in GridFM's architectural style
(the class \texttt{GridFMHeteroSurrogate}); nothing imports the released
\texttt{gridfm\_graphkit} model and no GridFM checkpoint is loaded,
which is why it appears only as ``scratch'' and has no zero-shot row. Comparisons
involving GridFM in Sections~5--8 are therefore against a from-scratch model of
our own. Section~\ref{sec:graphkit} repeats the OPFData experiments with the
released \texttt{gridfm\_graphkit} model to close that gap; the conclusions do
not change, but the reason to prefer one model over another does.

The pipelines:

\begin{itemize}\itemsep2pt
\item the \textbf{pandapower pipeline}, which generates scenarios locally with
\texttt{pandapower.runopp} and serves them through the branch-row direct-SI
parquet used for the earlier power-flow study (four grids: case14, case118,
case300, SimBench);
\item the \textbf{OPFData pipeline}, which consumes DeepMind's published
OPFData release (\texttt{arXiv:2406.07234}) natively through PyTorch
Geometric's \texttt{OPFDataset} (two grids: case14, case118).
\end{itemize}

The two pipelines share nothing but the task definition, the loss, and the
metric code. They differ in solver, perturbation design, per-unit convention and
file format (Section~\ref{sec:pipelines}). Running the same three models through
both is therefore a test of whether the model ranking is a property of the
models or of the dataset --- which turns out to be the most informative result
here (Section~\ref{sec:cross}).

Hyperparameters are matched across pipelines and objectives: 40 epochs, LR
$5\times10^{-4}$ and batch 8 for GridFM, LR $10^{-4}$ and batch 4 for LUMINA and
GridSFM. Training-set sizes are comparable (11{,}998--12{,}053 for pandapower,
13{,}500 for OPFData), so neither pipeline enjoys a data-volume advantage.

\paragraph{The two objectives.} \textbf{Shared} is a single voltage-only loss
applied to all three models ($w_{\rm mse}=1.0$, $w_{\rm phys}=10^{-2}$, limit
weight 1.0, band weight 0). \textbf{Native} trains each model with the loss from
its own repository (\texttt{native\_loss.py}). Evaluation is identical under
both --- the same \texttt{opf\_task.opf\_metrics} scores every run --- so results
are comparable across objectives even though the loss values are not.
The native objectives are:

\begin{itemize}\itemsep2pt
\item \textbf{GridSFM}: \texttt{gridsfm.loss.compute\_loss}, ten weighted terms
--- tanh-capped $\theta$/V/$P_g$/$Q_g$, BCE on feasibility, log-MSE on cost,
log1p KCL $P$/$Q$, branch flow $P$/$Q$, thermal loading and a thermal-limit
barrier. OPFData supplies \texttt{edge\_label} $(P_{ij},Q_{ij},P_{ji},Q_{ji})$
and \texttt{rate\_a}, so every term is live except the two feasibility terms,
which have no signal because OPFData ships solved instances only.
\item \textbf{LUMINA}: \texttt{lumina.model.opf.losses.ACOPFLossFunction},
weighting bus \emph{and} generator regression equally.
\item \textbf{GridFM}: \texttt{gridfm\_graphkit}'s
\texttt{MaskedReconstructionMSE} over $[V_M, V_A, P_G, Q_G, P_D, Q_D]$. Their
\texttt{PBELoss} is not usable here (it needs a six-column YFF/YFT edge
attribute that the OPFData adapter does not carry), so our physics term stands
in for it; this substitution is printed in every run log.
\end{itemize}

\paragraph{Task.} The model receives the loads and the decision space
(per-bus generation limits, cost slope, voltage band) and must produce the OPF
operating point. A voltage-only head is \emph{sufficient in principle}: given
the predicted profile the dispatch follows from \(S = V \odot \overline{YV}\),
so all three surrogates keep their existing per-bus heads.

\paragraph{Supervision, and why it is not neutral.}
\label{par:supervision}
The loss used throughout is
\[
\mathcal{L} = w_{\rm mse}\,\mathrm{MSE}(V_{\rm pred}, V_{\rm target})
 + w_{\rm phys}\,[\text{balance at free buses} + \text{limit violation at controllable buses}],
\]
so \textbf{bus voltages are the only supervised quantity}. Both pipelines also
carry the optimal dispatch, the objective value, and (pandapower only) the LMP
duals; none of these entered the loss. Dispatch and cost appear only in the
metrics, and the duals not at all.

This choice keeps the three models comparable, but it is \emph{not} architecture
neutral, and the results below must be read with that in mind:

\begin{itemize}\itemsep2pt
\item \textbf{GridFM} predicts voltages directly. The objective matches its
architecture exactly, and it is unaffected.
\item \textbf{LUMINA} also emits a generator head (\texttt{out\_dict.generator},
258 parameters, trained on ACOPF dispatch in the released weights) which this
setup never reads. Its bus head does not depend on it, so its voltage accuracy
is unaffected; only its dispatch capability goes unmeasured.
\item \textbf{GridSFM} is the one whose architecture this objective clearly does
not fit. Its angle prediction is built on a DC prior computed \emph{from its own
predicted} \(P_g\) (\texttt{model.py:234,243}) through a doubly detached path
(\texttt{dc\_prior.py:102,243}), so \texttt{head\_Pg} receives no gradient ---
confirmed by a single-batch probe. It also consumes the generation box
structurally, building its angle prior from the box midpoint --- which is what
made it uniquely sensitive to the conversion bug described in
Section~\ref{sec:gridsfm}. This supervision mismatch is a real limitation of the
comparison, but it was never the explanation of the observed failure.
\end{itemize}

The shared comparison is therefore a fair test of \emph{these models under
voltage-only supervision}, which flatters GridFM and penalises GridSFM. That was
stated as a caveat in earlier revisions; it is now measured
(Section~\ref{sec:native}), and the effect is large enough to reverse the
ranking on case118.

\section{The two pipelines}
\label{sec:pipelines}

\begin{table}[H]
\centering
\small
\begin{tabular}{p{0.20\textwidth}p{0.36\textwidth}p{0.36\textwidth}}
\toprule
 & \textbf{pandapower pipeline} & \textbf{OPFData pipeline} \\
\midrule
Source & generated locally, \texttt{pp.runopp} & DeepMind OPFData release 1 \\
Solver & pandapower interior point (PIPS) & PowerModels.jl + IPOPT \\
Perturbation & voltage setpoints 0.90--1.10, load scaling 0.60--1.40 (case300: 0.90--1.10) & per-load scaling with independent noise, as published \\
Topology & fixed per grid & fixed (perturbations \emph{not} enabled) \\
Format & branch-row direct-SI parquet, 20 rows/group & PyG \texttt{OPFDataset}, \texttt{HeteroData} \\
Per unit & rebased to 100~MVA, complex128 & already pu, \(S_{\rm base}=1\) \\
Grids used & case14, case118, case300, SimBench & case14, case118 \\
Train / val / test & \(\approx\)12k / 12k / 12k & 13{,}500 / 750 / 750 \\
Loader & \texttt{Dataset\_optimized\_}\-\texttt{complex\_columns}, lazy row groups & \texttt{opfdata\_pipeline.py}, PyG \texttt{Batch} \\
\bottomrule
\end{tabular}
\caption{The pipelines are independent implementations end to end. Both feed
the same three models through the same loss and the same metric code
(\texttt{opf\_task.py}).}
\end{table}

\paragraph{The decision spaces agree.} A useful consistency check fell out of
the metric masks. For a given grid the two pipelines independently identify the
same number of free and controllable buses per graph: case14 gives 9 free and 5
controllable in both, case118 gives 64 free and 54 controllable in both. The
two pipelines therefore agree on what the OPF problem \emph{is}, which is what
makes the cross-pipeline comparison in Section~\ref{sec:cross} meaningful.

\paragraph{OPFData cost-model validation.} The cost model was checked against
the dataset's own stored objective: recomputing the IPOPT objective from the
stored dispatch reproduces the published value to a ratio of 1.000000, and
recovering the dispatch from the stored voltages matches
\texttt{generator.y} to \(5.5\times10^{-7}\). The OPFData cost gaps below are
therefore measured against a verified reference.

\section{Datasets}

\begin{table}[H]
\centering
\small
\begin{tabular}{llrrrrl}
\toprule
Pipeline & Dataset & Buses & Scenarios & Dispatch buses & Target residual \\
\midrule
pandapower & case14   &  14 & 36{,}000 &  5 & \(6.8\times10^{-8}\) \\
           & case118  & 118 & 36{,}000 & 54 & \(1.6\times10^{-7}\) \\
           & case300  & 300 & 36{,}163 & 69 & \(3.2\times10^{-6}\) \\
           & SimBench &  94 & 36{,}000 & 70 & \(7.1\times10^{-7}\) (median) \\
\midrule
OPFData    & case14   &  14 & 15{,}000 &  5 & --- (verified vs.\ stored objective) \\
           & case118  & 118 & 15{,}000 & 54 & --- (verified vs.\ stored objective) \\
\bottomrule
\end{tabular}
\caption{The target residual is \(\max|\Delta P|\) in pu when the stored solution
is substituted into the dataset's own \(Y_{\rm bus}\); it is the noise floor
below which a model cannot be credited.}
\end{table}

Four preparation issues were found and corrected in the pandapower pipeline;
each would have silently produced wrong results. The OPFData pipeline, being
consumed as published, needed none of them --- which is part of its value as an
independent check.

\begin{itemize}\itemsep2pt
\item \textbf{Sign convention.} In the backbone files \texttt{S\_demand} is
already a signed injection, not positive demand. Reading it as demand put the
balance residual at \(1.07\) pu on the \emph{ground truth}; corrected it is
\(\sim\!10^{-7}\).
\item \textbf{Parquet row groups.} The sources use 1900--2000 rows per group.
Training shuffles, so each batch of 4 decoded a whole group: 3191~ms/batch
against 176~ms/batch at 20 rows per group, an 18\(\times\) penalty. All files
were re-laid-out.
\item \textbf{Angle datum.} pandapower's case118 places its slack at
\(+30^\circ\), so no bus sits near zero, while the other three grids use
\(0^\circ\). The datum is a gauge choice --- \(S = V\odot\overline{YV}\) is
invariant under a global phase rotation --- but it penalises any model that
cannot observe the bus angle state. LUMINA's bus schema has no such input, so it
paid \(\sim\!30^\circ\) before solving anything. Normalising the slack to
\(0^\circ\) left \(|V|\), \(S_{\rm start}\), \(S_{\rm newton}\) and the balance
residual bit-identical and moved only the angle metric
(LUMINA \(52.2^\circ \rightarrow 34.4^\circ\); GridSFM \(22.7^\circ \rightarrow
38.7^\circ\), i.e.\ the artifact had been flattering GridSFM).
\item \textbf{Generation limits.} The stock pandapower limits are violated by the
stored solutions, so the generating OPF used different ones. Limits are instead
the observed per-bus dispatch envelope widened by 10\%, which makes the
reference feasible by construction; the box is therefore the observed range, not
the true one. Balance and voltage columns do not depend on this choice.
\end{itemize}

\section{Metrics, and one that does not survive scrutiny}
\label{sec:metrics}

What a metric can mean differs per bus. Where no controllable generation sits,
the injection is known and the power-balance residual is a genuine error; where
the dispatch is free, the meaningful quantity is how far the implied generation
leaves its box. A plain residual over all buses would charge \emph{correct
redispatch} as error, so everything is masked by the controllable set
(\texttt{opf\_task.py}). Reported columns: voltage RMSE, angle RMSE, balance
\(\Delta P_\infty\) at free buses, generation-limit violation at controllable
buses, and the dispatch cost gap against the reference optimum.

\paragraph{The pandapower cost gaps are computed against a malformed generation
box.} The original conversion built the generation envelope on a redispatch
delta rather than on generation, producing a box centred at \(-28\)~pu on
case118. The full diagnosis is in Section~\ref{sec:gridsfm}; what matters here is
that the cost gap is evaluated against that box, so the original cost column is
not measuring what it claims to. Re-running on the corrected files
(\texttt{-{}-gen\_box injection}) leaves voltage accuracy essentially unchanged
but moves the cost gap by hundreds of points:

\begin{table}[H]
\centering
\small
\begin{tabular}{llrrrr}
\toprule
Grid & Run & \(|V|\) RMSE (orig.) & \(|V|\) RMSE (rev.) & cost gap (orig.) & cost gap (rev.) \\
\midrule
case14  & GridFM scratch  & 8.94e-4 & 9.90e-4 & --- & 0.89\% \\
case118 & GridFM scratch  & 1.32e-3 & 1.11e-3 & \(+3.20\%\) & \(-277.36\%\) \\
        & LUMINA scratch  & 8.81e-3 & 9.15e-3 & \(+8.14\%\) & \(+3391.19\%\) \\
        & LUMINA pre-FT   & 1.26e-2 & 1.32e-2 & \(+11.65\%\) & \(+149.78\%\) \\
case300 & GridFM scratch  & 6.57e-3 & 7.07e-3 & \(+156.44\%\) & \(-0.71\%\) \\
        & LUMINA scratch  & 2.02e-2 & 1.98e-2 & \(+749.21\%\) & \(+10.20\%\) \\
\bottomrule
\end{tabular}
\caption{Same models, same scenarios, near-identical voltage predictions --- and
cost gaps that move by hundreds of percentage points and change sign once the
generation box is built correctly. The voltage columns move by at most 13\%,
which is why the model ranking is unaffected by the correction.}
\end{table}

The voltage and balance metrics are stable across the correction; the cost gap is
not. \textbf{All model rankings in this report are therefore made on voltage and
balance metrics}, and the pandapower cost column should be disregarded entirely.
Even after correction the case118 cost gaps remain implausible
(\(-277\%\), \(+3391\%\)), which is consistent with the further finding that
case118's load cannot be recovered from its stored base generation --- so no
absolute cost reference exists for that grid. The OPFData cost gaps are the only
ones in this report measured against a verified reference
(Section~\ref{sec:pipelines}). A negative gap means the dispatch is cheaper than
optimal, i.e.\ infeasible.

\section{Results: pandapower pipeline}

\begin{table}[H]
\centering
\small
\begin{tabular}{llrrrr}
\toprule
Grid & Run & RMSE & \(|V|\) RMSE & \(\theta\) [deg] & \(\Delta P_\infty\) (unmasked) \\
\midrule
case14 & \textbf{GridFM scratch} & \textbf{1.007e-3} & \textbf{8.94e-4} & \textbf{0.027} & 3.09e-1 \\
       & LUMINA pre-FT          & 8.225e-3 & 3.47e-3 & 0.427 & 4.05e-1 \\
       & LUMINA scratch         & 1.074e-2 & 3.61e-3 & 0.579 & 4.20e-1 \\
       & GridSFM scratch$^{\dagger}$ & 3.795e-2 & 2.29e-3 & 2.171 & 4.44e-1 \\
       & GridSFM pre-FT$^{\dagger}$  & 5.263e-2 & 3.38e-3 & 3.009 & 8.87e-1 \\
\midrule
SimBench & \textbf{GridFM scratch} & \textbf{7.594e-3} & \textbf{4.74e-3} & \textbf{0.340} & 1.87 \\
         & LUMINA scratch         & 1.216e-1 & 4.62e-2 & 6.445 & 3.38 \\
         & LUMINA pre-FT          & 1.217e-1 & 4.75e-2 & 6.422 & 2.93 \\
         & GridSFM scratch$^{\dagger}$ & 2.489e-1 & 5.57e-2 & 13.900 & 1.05 \\
         & GridSFM pre-FT$^{\dagger}$  & 4.582e-1 & 5.57e-2 & 26.061 & 2.27 \\
\midrule
case118 & \textbf{GridFM scratch} & \textbf{6.459e-3} & \textbf{1.32e-3} & \textbf{0.362} & 1.53e-1$^{\ast}$ \\
        & LUMINA scratch         & 2.445e-2 & 8.81e-3 & 1.307 & 1.00$^{\ast}$ \\
        & LUMINA pre-FT          & 3.023e-2 & 1.26e-2 & 1.574 & 9.37e-1$^{\ast}$ \\
        & GridSFM$^{\dagger\ddagger}$ & \multicolumn{4}{l}{frozen; did not complete (epochs 25 / 28 of 40)} \\
\midrule
case300 & \textbf{GridFM scratch} & \textbf{2.847e-2} & \textbf{6.57e-3} & \textbf{1.587} & 1.65$^{\ast}$ \\
        & LUMINA pre-FT          & 7.144e-2 & 1.99e-2 & 3.931 & 6.49$^{\ast}$ \\
        & LUMINA scratch         & 7.132e-2 & 2.02e-2 & 3.919 & 1.15e1$^{\ast}$ \\
        & GridSFM$^{\dagger\ddagger}$ & \multicolumn{4}{l}{frozen; did not complete (epochs 38 / 39 of 40)} \\
\bottomrule
\end{tabular}
\caption{Final test results.
$^{\ast}$masked balance at free buses; elsewhere the logs record only the
unmasked residual, so the columns are not comparable across grid blocks.
$^{\dagger}$\textbf{every} GridSFM row here --- including the case14 rows, which
did not freeze --- was produced on the defective conversion
(Section~\ref{sec:gridsfm}) and measures the bug, not the model. Not freezing is
not the same as being valid: case14 trained, but against a generation box still
built on a redispatch delta, so its RMSE reflects a decision space that does not
describe the OPF problem. No GridSFM number in this table should be quoted as a
capability measurement, and the corrected case14 numbers are not yet available
to replace them.
$^{\ddagger}$still running at the time of writing, both in the frozen regime, so
they cannot change any conclusion.}
\end{table}

\begin{table}[H]
\centering
\small
\begin{tabular}{llrrr}
\toprule
Grid & Run & RMSE & \(|V|\) RMSE & \(\theta\) [deg] \\
\midrule
case14  & GridFM scratch  & \textbf{1.156e-3} & \textbf{9.90e-4} & \textbf{0.034} \\
        & LUMINA scratch  & 6.739e-3 & 2.81e-3 & 0.351 \\
        & LUMINA pre-FT   & 1.205e-2 & 4.23e-3 & 0.646 \\
        & GridSFM scratch & 4.599e-2 & 4.79e-3 & 2.621 \\
        & GridSFM pre-FT  & 4.860e-2 & 2.15e-3 & 2.782 \\
\midrule
case118 & GridFM scratch  & \textbf{6.342e-3} & 1.11e-3 & \textbf{0.358} \\
        & LUMINA pre-FT   & 3.059e-2 & 1.32e-2 & 1.581 \\
        & LUMINA scratch  & 3.133e-2 & 9.15e-3 & 1.717 \\
        & GridSFM pre-FT  & 5.065e-2 & \textbf{2.07e-3} & 2.900 \\
        & GridSFM scratch$^{\S}$ & 3.546e-1 & 1.02e-2 & 20.306 \\
\midrule
case300 & GridFM scratch  & \textbf{2.973e-2} & \textbf{7.07e-3} & \textbf{1.655} \\
        & LUMINA scratch  & 6.992e-2 & 1.98e-2 & 3.843 \\
        & LUMINA pre-FT   & 7.276e-2 & 2.01e-2 & 4.006 \\
        & GridSFM pre-FT$^{\S}$  & 1.649e-1 & 1.76e-2 & 9.396 \\
        & GridSFM scratch$^{\S}$ & 3.166e-1 & 1.63e-2 & 18.118 \\
\bottomrule
\end{tabular}
\caption{pandapower pipeline on the \emph{corrected} conversion
(\texttt{-{}-gen\_box injection}), voltage-only objective. These supersede the
GridSFM rows of Table~4. GridSFM no longer freezes but remains last on every
grid: correcting the bug made the numbers valid, it did not change the ranking.
Its \(|V|\) is a different story --- 2.07e-3 on case118, four times better than
LUMINA and close to GridFM --- and only its angle is poor, the same
magnitude-good/angle-bad signature that the native objective repairs
(Section~\ref{sec:native}). $^{\S}$scored from the epoch-37/38 checkpoint after
the run exceeded the 24\,h limit before its own test pass, so these three are
marginally pessimistic relative to a full 40 epochs. They do not change the
picture: on corrected data GridSFM's angle error is still 9--20\(^\circ\) on
case118 scratch and case300, against \(0.36\)--\(1.66^\circ\) for GridFM.
Correcting the conversion removed the freeze; it did not make GridSFM
competitive under a voltage-only objective.}
\end{table}

\section{Results: OPFData pipeline, shared objective}

\begin{table}[H]
\centering
\small
\begin{tabular}{llrrrrrr}
\toprule
Grid & Run & RMSE & \(|V|\) RMSE & \(\theta\) [deg] & \(\Delta P_\infty\) & limit \(P\) & cost gap \\
\midrule
case14 & \textbf{GridFM scratch} & \textbf{5.950e-4} & \textbf{4.52e-4} & \textbf{0.022} & \textbf{8.07e-3} & \textbf{4.47e-3} & \textbf{0.14\%} \\
       & LUMINA pre-FT          & 7.683e-3 & 2.03e-3 & 0.425 & 7.70e-2 & 4.95e-2 & 13.55\% \\
       & LUMINA scratch         & 7.771e-3 & 1.77e-3 & 0.434 & 7.10e-2 & 6.30e-2 & 12.69\% \\
       & GridSFM pre-FT         & 7.919e-3 & 1.57e-3 & 0.445 & 3.97e-2 & 3.93e-2 & 7.24\% \\
       & GridSFM scratch$^{\dagger}$ & 3.105e-2 & 3.03e-2 & 0.380 & 5.82e-2 & 1.20e-1 & 30.00\% \\
\midrule
case118 & \textbf{GridFM scratch} & \textbf{2.271e-2} & \textbf{3.32e-3} & \textbf{1.287} & \textbf{1.15e-1} & \textbf{1.47e-1} & 0.86\% \\
        & LUMINA scratch         & 3.288e-2 & 1.19e-2 & 1.755 & 1.13 & 1.90 & \(-1.71\%\) \\
        & LUMINA pre-FT          & 3.694e-2 & 1.46e-2 & 1.944 & 1.69 & 1.80 & 2.79\% \\
        & GridSFM scratch$^{\dagger}$ & 1.196e-1 & 3.39e-2 & 6.571 & 3.93e-1 & 2.36e-1 & 21.21\% \\
        & GridSFM pre-FT$^{\dagger}$  & 3.507e-1 & 3.39e-2 & 20.000 & 1.71 & 5.85 & 11.53\% \\
\bottomrule
\end{tabular}
\caption{Voltage-only objective, all runs complete (40 epochs). The
\(\Delta P_\infty\) and limit columns here are batch-averaged and therefore
understate the true worst bus --- see Section~\ref{sec:resid}; pooled residuals
are given there for the native runs only.
$^{\dagger}$degenerate: three of these report a \(|V|\) RMSE of exactly
3.3899e-02 or 3.03e-02 with a flat loss, the collapsed-magnitude signature of
Section~\ref{sec:gridsfm}. Under GridSFM's own objective none of them collapses
(Section~\ref{sec:native}).}
\end{table}

\section{Results: OPFData pipeline, native objectives}
\label{sec:native}

Each model trained with the loss from its own repository, scored with the same
metrics. Every configuration was run independently on two clusters (alex A40,
helma H100) as a reproducibility check.

\begin{table}[H]
\centering
\small
\begin{tabular}{llrrrrr}
\toprule
Grid & Run & RMSE (helma) & RMSE (alex) & \(|V|\) RMSE & \(\theta\) [deg] & cost gap \\
\midrule
case14 & GridFM scratch  & 1.040e-3 & 4.300e-3 & 4.97e-4 & 0.052 & 1.48\% \\
       & GridSFM scratch & 4.710e-3 & 1.918e-3 & 1.18e-3 & 0.261 & 3.08\% \\
       & GridSFM pre-FT  & 5.327e-3 & 5.016e-3 & 1.32e-3 & 0.296 & 3.40\% \\
       & LUMINA pre-FT   & 7.247e-3 & 7.468e-3 & 1.80e-3 & 0.402 & \(-6.28\%\) \\
       & LUMINA scratch  & 7.523e-3 & 7.189e-3 & 1.71e-3 & 0.420 & \(-12.24\%\) \\
\midrule
case118 & \textbf{GridSFM pre-FT} & \textbf{2.080e-3} & \textbf{2.092e-3} & \textbf{7.53e-4} & \textbf{0.111} & 2.69\% \\
        & GridFM scratch  & 2.321e-2 & 2.311e-2 & 4.15e-3 & 1.308 & 4.48\% \\
        & LUMINA scratch  & 2.637e-2 & 2.693e-2 & 5.04e-3 & 1.483 & 19.62\% \\
        & LUMINA pre-FT   & 3.143e-2 & 3.134e-2 & 5.64e-3 & 1.772 & \(-2.94\%\) \\
        & GridSFM scratch$^{\ast}$ & 3.657e-2 & \textbf{1.809e-3} & 3.39e-2 & 0.786 & \(-3.60\%\) \\
\bottomrule
\end{tabular}
\caption{Native objectives. \(|V|\), \(\theta\) and cost columns are the helma
run. $^{\ast}$the two clusters disagree by \(20\times\) on this row: the helma
run collapsed (\(|V|\) pinned at 3.39e-2) while the alex run reached 1.809e-3
(\(|V|\) 5.33e-4, \(\theta\) 0.099\(^\circ\)) --- the best result anywhere in
this study. GridSFM-from-scratch therefore collapses \emph{stochastically}; its
pretrained runs did not collapse on either cluster.}
\end{table}

\paragraph{The objective, not the architecture, decided the earlier ranking.}
On case118 GridSFM goes from last under the shared objective to first under its
own, and the margin is not marginal:

\begin{table}[H]
\centering
\small
\begin{tabular}{llrrr}
\toprule
Grid & Run & shared & native & factor \\
\midrule
case118 & GridSFM pre-FT  & 3.507e-1 & 2.080e-3 & \textbf{169\(\times\)} \\
        & GridSFM scratch & 1.196e-1 & 3.657e-2 / 1.809e-3 & 3.3--66\(\times\) \\
        & LUMINA scratch  & 3.288e-2 & 2.637e-2 & 1.25\(\times\) \\
        & LUMINA pre-FT   & 3.694e-2 & 3.143e-2 & 1.18\(\times\) \\
        & GridFM scratch  & 2.271e-2 & 2.321e-2 & 0.98\(\times\) (worse) \\
\midrule
case14  & GridSFM scratch & 3.105e-2 & 4.710e-3 & 6.6\(\times\) \\
        & GridSFM pre-FT  & 7.919e-3 & 5.327e-3 & 1.5\(\times\) \\
        & LUMINA pre-FT   & 7.683e-3 & 7.247e-3 & 1.06\(\times\) \\
        & GridFM scratch  & 5.950e-4 & 1.040e-3 & 0.57\(\times\) (worse) \\
\bottomrule
\end{tabular}
\caption{Test RMSE under each objective. The ordering of the \emph{gains} is the
result: GridSFM \(\gg\) LUMINA \(>\) GridFM, which is exactly the order in which
the shared voltage-only objective departs from what each architecture expects.}
\end{table}

\subsection*{Power-balance residuals}
\label{sec:resid}

Voltage RMSE says how close the predicted operating point is to the reference
one; it does not say whether that point satisfies the network equations. The
residual \(\Delta S = S_{\rm implied} - S_{\rm load}\) at buses with no
controllable generation does, and it is the quantity an operator would actually
care about. All statistics below are pooled over every free bus in the test set
and are for the native-objective runs.

\begin{table}[H]
\centering
\small
\begin{tabular}{llrrrrrr}
\toprule
Grid & Run & \(\Delta P^{\infty}\) & \(\Delta Q^{\infty}\) & \(\overline{|\Delta P|}\) & \(\overline{|\Delta Q|}\) & p95\(|\Delta P|\) & p95\(|\Delta Q|\) \\
\midrule
case14 & GridFM scratch  & \textbf{2.450e-2} & \textbf{2.166e-2} & 3.963e-3 & \textbf{1.854e-3} & \textbf{1.023e-2} & \textbf{5.053e-3} \\
       & GridSFM pre-FT  & 3.143e-2 & 3.888e-2 & \textbf{3.470e-3} & 3.934e-3 & 1.186e-2 & 1.134e-2 \\
       & GridSFM scratch & 4.110e-2 & 3.551e-2 & 3.462e-3 & 4.365e-3 & 1.200e-2 & 1.439e-2 \\
       & LUMINA pre-FT   & 1.484e-1 & 6.460e-2 & 2.998e-2 & 1.274e-2 & 8.574e-2 & 3.358e-2 \\
       & LUMINA scratch  & 1.739e-1 & 7.542e-2 & 3.957e-2 & 1.261e-2 & 1.140e-1 & 3.180e-2 \\
\midrule
case118 & \textbf{GridSFM pre-FT} & \textbf{2.646e-1} & 3.125e-1 & \textbf{9.011e-3} & \textbf{1.101e-2} & \textbf{3.566e-2} & \textbf{3.712e-2} \\
        & GridFM scratch  & 3.975e-1 & \textbf{2.376e-1} & 3.045e-2 & 1.830e-2 & 9.613e-2 & 5.603e-2 \\
        & GridSFM scratch & 6.411e-1 & 2.436e+0 & 4.519e-2 & 3.129e-1 & 1.209e-1 & 1.717e+0 \\
        & LUMINA pre-FT   & 5.297e+0 & 1.331e+0 & 1.862e-1 & 8.455e-2 & 5.989e-1 & 2.878e-1 \\
        & LUMINA scratch  & 6.477e+0 & 3.692e+0 & 1.569e-1 & 1.000e-1 & 6.653e-1 & 2.517e-1 \\
\bottomrule
\end{tabular}
\caption{Power-balance residuals in pu on a 100~MVA base (so 1e-2~pu \(=\)
1~MW/MVAr), pooled over all free buses of the test set. Rows are ordered by
\(\overline{|\Delta P|}\) within each grid.}
\end{table}

\paragraph{The residual view mostly agrees with RMSE, and where it does not it
is more informative.} On case118 GridSFM pre-FT leads on every column except
\(\Delta Q^{\infty}\), where GridFM is better (2.376e-1 vs 3.125e-1). So GridSFM
is the better surrogate on average --- \(3.4\times\) lower
\(\overline{|\Delta P|}\) --- while GridFM has the better \emph{worst-case}
reactive behaviour. An RMSE-only table cannot show that trade-off.

On case14 the picture is mixed rather than a GridFM sweep: GridSFM's
\(\overline{|\Delta P|}\) is marginally lower (3.46--3.47e-3 vs 3.96e-3) while
GridFM wins the other five columns. Given the scratch-run variance of
Section~\ref{sec:native}, case14 should be read as a tie between the two.

LUMINA is decisively last on both grids by every residual measure, and by a
wider margin than its voltage RMSE suggests: on case118 its worst-bus \(\Delta
P\) is 5.3--6.5~pu, i.e.\ 530--650~MW unbalanced at a single bus, against
0.26--0.40~pu for the other two. Its bus voltages are only \(1.4\times\) worse
than GridFM's in RMSE but its residuals are \(6\)--\(16\times\) worse. Accurate
voltages do not imply a physically consistent operating point, which is the main
argument for reporting residuals alongside RMSE.

\paragraph{A correction to how \(\Delta P^{\infty}\) was computed.} The
per-epoch metric block accumulates every quantity as a batch-size-weighted mean
across batches. That is correct for means but not for a maximum: the
\(\Delta P^{\infty}\) reported in earlier revisions (and in Table~6) is the
\emph{average of per-batch maxima}, which understates the true worst bus, and a
percentile cannot be recovered from per-batch summaries at all. The table above
pools the per-bus residuals and summarises once. The difference is not cosmetic
--- \(3\)--\(6\times\) on the runs where both are available:

\begin{center}
\small
\begin{tabular}{lrr}
\toprule
Run & \(\Delta P^{\infty}\) as reported & \(\Delta P^{\infty}\) pooled \\
\midrule
case118 GridFM scratch  & 1.15e-1 & 3.975e-1 \\
case118 GridSFM pre-FT  & 6.72e-2 & 2.646e-1 \\
case118 LUMINA scratch  & 1.13    & 6.477e+0 \\
case14  GridFM scratch  & 8.07e-3 & 2.450e-2 \\
\bottomrule
\end{tabular}
\end{center}

The pandapower driver already pooled correctly, so its \(\Delta P^{\infty}\)
column was never affected --- but that also means the two pipelines were
reporting the same symbol under two different definitions, and cross-pipeline
\(\Delta P^{\infty}\) comparisons in earlier revisions should be discarded.

The mechanism is visible in the angle column. Under the shared objective
GridSFM's case118 angle error is \(20.000^\circ\) (a clip bound); under its own
objective it is \(0.111^\circ\), twelve times better than GridFM's
\(1.308^\circ\). GridSFM builds its angle prior from its own predicted \(P_g\)
through a doubly detached path, so a voltage-only loss leaves \texttt{head\_Pg}
untrained and the angle unanchored. Supervising \(P_g\) --- which its native
loss does --- repairs it. This was stated as a hypothesis in earlier revisions
and is now confirmed.

GridFM is the mirror image: it gets \emph{worse} under the native objective on
both grids, because the shared voltage-only loss was already exactly what its
architecture emits, and adding \(P_G\)/\(Q_G\)/\(P_D\)/\(Q_D\) reconstruction
targets dilutes it. LUMINA sits between the two, gaining 6--25\%, consistent
with its generator head being trained rather than ignored.

\paragraph{Pretrained runs replicate; scratch runs do not.} The two clusters
agree to 0.3--2\% on every \emph{pretrained} case118 row --- GridSFM pre-FT
2.080e-3 vs 2.092e-3, GridFM 2.321e-2 vs 2.311e-2, LUMINA pre-FT 3.143e-2 vs
3.134e-2 --- so the \(11\times\) case118 result is solid.

Scratch runs are a different matter and the variance is not small:

\begin{itemize}\itemsep2pt
\item GridFM case14: 1.040e-3 (helma) vs 4.300e-3 (alex), \(4.1\times\).
\item GridSFM case14: 4.710e-3 vs 1.918e-3, \(2.5\times\) --- which swaps its
rank against GridFM.
\item GridSFM case118: 3.657e-2 vs 1.809e-3, \(20\times\). The helma run
collapsed its magnitude head; the alex run did not, and produced the best number
in this study.
\end{itemize}

The last is the important one: \textbf{the magnitude collapse is stochastic in
scratch runs.} It is not a property of the model, the objective or the pipeline
--- the same configuration collapsed on one cluster and reached 1.809e-3 on the
other. Every collapsed run in this document is a scratch run; no pretrained run
collapsed anywhere. Pretraining appears to place GridSFM outside the basin that
leads to collapse, which is a concrete reason to prefer its released checkpoint
that has nothing to do with transfer of grid knowledge.

Consequently \textbf{no ordering should be quoted from a single scratch run},
on either grid. Seed repeats are required before any scratch row is used to rank
models. Pretrained rows and case118 margins are safe.

\subsection*{The released GridFM model versus our re-implementation}
\label{sec:graphkit}

Every GridFM number above comes from a local model written in GridFM's style
(Section~1). LUMINA and GridSFM are run from their own repositories, so this was
the one place where the comparison was not like-for-like. The released model,
\texttt{gridfm\_graphkit.models.gnn\_heterogeneous\_gns.GNS\_heterogeneous}, was
therefore run on the same two grids under both objectives, at the same capacity
(hidden 48, 12 layers, 8 heads --- the released config, which the mirror already
matched).

The released model is not a renamed copy. It injects a physics correction at
\emph{every} layer (branch flows \(\rightarrow\) node injections \(\rightarrow\)
\texttt{physics\_mlp}), bounds \(|V|\) with a sigmoid against the per-bus band
rather than a fixed clamp, carries a generator head, pins known quantities to
their inputs through \texttt{mask\_dict}, and --- most relevant here --- ends in
a task-specific physics decoder that recovers \(Q_g\) at PV and slack buses from
the nodal reactive balance instead of learning it.

\begin{table}[H]
\centering
\small
\begin{tabular}{lllrrrrr}
\toprule
Grid & Obj. & Implementation & RMSE & \(|V|\) RMSE & \(\Delta P^{\infty}\) & \(\overline{|\Delta P|}\) & p95\(|\Delta P|\) \\
\midrule
case14 & shared & mirror   & \textbf{5.950e-4} & \textbf{4.52e-4} & 2.450e-2 & 3.963e-3 & 1.023e-2 \\
       &        & released & 6.349e-4 & 5.00e-4 & \textbf{1.123e-2} & \textbf{1.245e-3} & \textbf{4.039e-3} \\
       & native & mirror   & \textbf{1.040e-3} & \textbf{4.97e-4} & \textbf{2.450e-2}$^{\ast}$ & \textbf{3.963e-3}$^{\ast}$ & \textbf{1.023e-2}$^{\ast}$ \\
       &        & released & 1.212e-3 & 7.46e-4 & 7.719e-2 & 8.608e-3 & 3.292e-2 \\
\midrule
case118 & shared & mirror   & \textbf{2.271e-2} & \textbf{3.32e-3} & 3.975e-1 & 3.045e-2 & 9.613e-2 \\
        &        & released & 2.374e-2 & 4.74e-3 & \textbf{2.673e-1} & \textbf{1.640e-2} & \textbf{4.504e-2} \\
        & native & mirror   & \textbf{2.321e-2} & \textbf{4.15e-3} & \textbf{3.975e-1} & \textbf{3.045e-2} & \textbf{9.613e-2} \\
        &        & released & 2.904e-2 & 1.43e-2 & 4.969e+0 & 2.364e-1 & 7.219e-1 \\
\bottomrule
\end{tabular}
\caption{Local re-implementation versus the released \texttt{gridfm\_graphkit}
model, same capacity and same data. $^{\ast}$the mirror's residual columns were
pooled from its native-objective checkpoint (Section~\ref{sec:resid}); its
shared-objective residuals were not re-scored, so the two mirror rows repeat the
same figures and only the RMSE columns differ between them.}
\end{table}

\paragraph{The released model trades a little voltage accuracy for a much more
consistent operating point.} Under the shared objective its RMSE is
\emph{worse} by 5--7\% on both grids, while its power-balance residuals are far
better: \(\overline{|\Delta P|}\) improves \(3.2\times\) on case14 (1.245e-3 vs
3.963e-3) and \(1.9\times\) on case118 (1.640e-2 vs 3.045e-2), with
\(\Delta P^{\infty}\) improving \(2.2\times\) and \(1.5\times\). That is exactly
the behaviour its architecture would predict: the per-layer physics feedback and
the physics decoder push the prediction towards satisfying the network equations
rather than towards matching the reference voltages.

This is a concrete instance of the point made in Section~\ref{sec:resid}:
ranked on RMSE the mirror wins, ranked on residuals the released model wins, and
the two answers are both correct about different things. \textbf{Which model is
``better'' here depends on whether the downstream use needs a voltage estimate
or a feasible operating point.}

\paragraph{GridFM's own loss hurts the released model too, and more.} Under the
native objective the released model degrades sharply on case118 ---
\(\overline{|\Delta P|}\) rises from 1.640e-2 to 2.364e-1, a factor of 14 --- and
its \(|V|\) RMSE triples. The mirror showed the same direction
(Section~\ref{sec:native}) but not the same magnitude. The likely reason is
mechanical: with \texttt{MaskedReconstructionMSE} the supervised \(P_G\) and
\(Q_G\) come out of the physics decoder, so an error signal on them now
propagates back through the physics path that otherwise regularises the voltage
prediction. This is a hypothesis; it has not been isolated.

\paragraph{What this does and does not change.} The GridSFM result of
Section~\ref{sec:native} survives: on case118 under native objectives GridSFM
reaches 2.080e-3 RMSE and \(\overline{|\Delta P|}\) 9.011e-3, still ahead of the
released GridFM's best case118 configuration (2.374e-2 and 1.640e-2) by
\(11\times\) and \(1.8\times\) respectively. Swapping in the real GridFM does
not overturn any ranking in this report; it does change the reason to prefer
one model over another, and it removes the last respect in which GridFM was
being evaluated differently from the other two.

\paragraph{One asymmetry remains.} \texttt{gridfm\_graphkit} ships no weights
--- no checkpoint file and no HuggingFace reference in the repository --- so
GridFM is scratch-only under either implementation, while LUMINA and GridSFM
have public checkpoints. No GridFM row in this report is, or can be, a
pretrained one.

\paragraph{Conversion correctness.} The released model computes branch flows
itself and so needs the per-branch admittance split (\(Y_{ff}\), \(Y_{ft}\))
across ten edge columns, where this project carries only the two-column Y-bus
off-diagonal. The split was rebuilt from the branch model and checked by
re-assembling a Y-bus from it and comparing against the dataset's own: relative
error \(1.8\times10^{-16}\) (case14) and \(7.3\times10^{-17}\) (case118) on
OPFData, and \(1.4\)--\(2.6\times10^{-16}\) across case118, case300, GBnetwork
and SimBench on the parquet pipeline. The check is not ceremonial: a first
version applied the per-unit bus bases a second time --- the parquet loader has
already applied them --- and missed the Y-bus by a relative \(7.5\times10^{2}\)
without raising anything.

\section{Cross-pipeline comparison}
\label{sec:cross}

Both pipelines cover case14 and case118 with matched hyperparameters and
comparable training-set sizes, so the same model can be compared head to head.

\begin{table}[H]
\centering
\small
\begin{tabular}{llrrr}
\toprule
Grid & Run & pandapower RMSE & OPFData RMSE & ratio \\
\midrule
case14 & GridFM scratch  & 1.007e-3 & 5.950e-4 & 0.59 \\
       & LUMINA pre-FT   & 8.225e-3 & 7.683e-3 & 0.93 \\
       & LUMINA scratch  & 1.074e-2 & 7.771e-3 & 0.72 \\
       & GridSFM pre-FT  & 4.860e-2$^{c}$ & 7.919e-3 & 0.16 \\
       & GridSFM scratch & 4.599e-2$^{c}$ & 3.105e-2 & 0.68 \\
\midrule
case118 & GridFM scratch & 6.459e-3 & 2.271e-2 & \textbf{3.52} \\
        & LUMINA scratch & 2.445e-2 & 3.288e-2 & 1.34 \\
        & LUMINA pre-FT  & 3.023e-2 & 3.694e-2 & 1.22 \\
\bottomrule
\end{tabular}
\label{tab:cross}
\caption{Ratio \(<1\) means the model does better on OPFData. For
GridFM and LUMINA the \emph{ordering} is preserved on both grids in both
pipelines, but the margins are not. $^{c}$GridSFM's pandapower entries are from
the \emph{corrected} conversion (Table~5), so these two ratios are valid; both
say GridSFM does markedly better on OPFData than on pandapower under the same
objective. No case118 GridSFM ratio is given because both of its OPFData
shared-objective runs are the collapsed ones.}
\end{table}

\paragraph{The GridFM--LUMINA ranking is a property of the models; the margin is
not.} GridFM leads LUMINA on both grids in both pipelines. That this survives a
complete change of solver, perturbation design, per-unit convention and file
format --- and survives the conversion correction, which moves their voltage
RMSE by at most 13\% --- is good evidence that this part of the ranking is real
and not a dataset artifact.

The \emph{size} of the lead is another matter. GridFM's margin over LUMINA is
\(8\times\) on pandapower case14 and \(13\times\) on OPFData case14, but on
case118 it is \(3.8\times\) on pandapower and only \(1.4\times\) on OPFData. Any
claim of the form ``GridFM is an order of magnitude better'' is pipeline
dependent and should not be made.

\paragraph{OPFData case118 is harder for GridFM specifically.} GridFM is
\(3.5\times\) worse on OPFData case118 than on pandapower case118, while LUMINA
is only \(1.2\)--\(1.3\times\) worse. Since the decision spaces are identical
(64 free / 54 controllable buses in both), the difference lies in the scenario
distribution: OPFData perturbs every load independently, whereas the pandapower
backbone applies a global load-scaling factor with voltage-setpoint variation.
The independent-perturbation design plausibly produces a higher-dimensional
target manifold that penalises the highest-capacity model's fit most --- but
this is a hypothesis suggested by the numbers, not something these runs
establish.

\paragraph{GridSFM is the pipeline- and objective-sensitive one.} Its ratios
above use the \emph{corrected} pandapower runs; the original-conversion numbers
would measure the bug rather than the model. Read across all three settings the
spread is stark. On case14 GridSFM goes 4.599e-2 (corrected pandapower,
shared objective) \(\rightarrow\) 3.105e-2 (OPFData, shared) \(\rightarrow\)
4.710e-3 (OPFData, native); on case118, 5.065e-2 \(\rightarrow\) 1.196e-1
\(\rightarrow\) 2.080e-3. That is a spread of up to two orders of magnitude for
one architecture on one grid, larger than any architecture-to-architecture
difference reported anywhere in this document. \textbf{What GridSFM is worth
cannot be stated without also stating the pipeline and the objective.}

\section{Zero-shot transfer}

\begin{table}[H]
\centering
\small
\begin{tabular}{lrrrrrr}
\toprule
 & \multicolumn{4}{c}{pandapower} & \multicolumn{2}{c}{OPFData} \\
\cmidrule(lr){2-5}\cmidrule(lr){6-7}
Released weights, no training & case14 & case118 & case300 & SimBench & case14 & case118 \\
\midrule
LUMINA \(\theta\) [deg]  & 12.87 & 34.42 & 21.83 & 38.96 & 10.03 & 12.69 \\
GridSFM \(\theta\) [deg] & 11.18 & 38.73 & 169.24 & 39.29 & 0.76 & 18.04 \\
LUMINA \(|V|\) RMSE      & 2.47e-2 & 3.74e-2 & 3.51e-2 & 6.15e-2 & 2.05e-2 & 2.14e-2 \\
GridSFM \(|V|\) RMSE     & 6.40e-2 & 3.67e-2 & 4.69e-2 & 1.02e-1 & 3.35e-2 & 3.54e-2 \\
\bottomrule
\end{tabular}
\caption{Neither model transfers reliably, but both do markedly better on
OPFData than on the pandapower grids. GridSFM's \(0.76^\circ\) zero-shot angle
on OPFData case14 is its single best result anywhere in this study.
Its \(169^\circ\) on pandapower case300 is near-maximal error (angles wrap at
\(180^\circ\)), i.e.\ that prediction carries no information.}
\end{table}

The zero-shot gap between pipelines is itself informative: both released models
were pretrained on pglib-derived data, which is what OPFData contains, and
neither was trained on locally generated pandapower snapshots or on SimBench's
CGMES-derived grid. GridSFM's angle error falls from \(11.2^\circ\) to
\(0.76^\circ\) on the same nominal grid (case14) purely by changing which
pipeline produced the scenarios.

\section{Findings}

\paragraph{There is no objective-independent winner.} Under the shared
voltage-only objective GridFM leads on every grid. Under each model's own
objective, GridSFM wins case118 by \(11\times\) and case14 is a statistical tie
between GridFM and GridSFM. The single most important result in this report is
that \textbf{the ranking is a property of the training objective, not of the
models}, and the earlier ``GridFM wins everywhere'' headline was an artifact of
choosing an objective that matches GridFM's output structure.

What survives objective changes is narrower: GridFM is the strongest surrogate
\emph{when only voltages are supervised}, and it is the least sensitive to which
objective is used (0.57--0.98\(\times\)). GridSFM is the strongest \emph{when
trained as its authors intended}, and the most sensitive by two orders of
magnitude. Which of these matters depends on whether a deployment can supply
dispatch, cost and flow targets --- OPFData can; the pandapower backbone
currently does not.

One caveat still attaches to GridFM specifically: it is the only model never
evaluated from a pretrained checkpoint, because none is published. The
re-implementation caveat has been discharged --- the released model was run on
the same grids and objectives (Section~\ref{sec:graphkit}) and changes no
ranking.

The qualification is not cosmetic. The training objective is voltage-only
(\S\ref{par:supervision}), which is precisely what GridFM's architecture emits,
while GridSFM's angle path depends on a generator head that this objective never
trains, and LUMINA's dispatch head goes unread. Part of the measured margin is
therefore attributable to the alignment between the loss and GridFM's output
structure rather than to modelling capacity. The ranking of GridFM against
LUMINA is nonetheless meaningful, since LUMINA's voltage head is independent of
its unused dispatch head; the ranking against GridSFM is not
(Section~\ref{sec:gridsfm}).

GridFM also carries the largest parameter count: 20.0~M, against GridSFM's
15.1~M and LUMINA's 2.4~M, and runs at batch 8 where the other two run at batch
4. Neither has been controlled for.

\paragraph{Grid size is the binding constraint under the shared objective.} On the
pandapower pipeline GridFM's voltage RMSE runs 1.01e-3 \(\rightarrow\) 6.46e-3
\(\rightarrow\) 2.85e-2 as the grid goes 14 \(\rightarrow\) 118
\(\rightarrow\) 300 buses, a \(28\times\) degradation, and every other model
degrades in step. At 300 buses no model is deployable. This is the result with
the most practical weight and it is independent of model choice.

\paragraph{LUMINA's pretrained weights give no reliable advantage.} This
replicates cleanly across both pipelines and is the most robust negative result
here. On pandapower, fine-tuning beats scratch on case14 (8.23e-3 vs 1.07e-2)
but scratch wins on case118 (2.44e-2 vs 3.02e-2) and SimBench is a tie
(1.2163e-1 vs 1.2173e-1). On OPFData the two are within 1\% on case14 (7.68e-3
vs 7.77e-3) and scratch is \emph{ahead} on case118 (3.29e-2 vs 3.69e-2).
Averaged over six grid/pipeline combinations the released weights are worth
nothing once training is allowed. This is not because the checkpoint is already
optimal --- its zero-shot errors are large --- but because 40 epochs of
in-distribution training erases whatever it provides.

\paragraph{LUMINA's magnitude transfers, its angle does not.} Zero-shot
\(|V|\) stays within 2.1--6.2e-2 across all six combinations while \(\theta\)
ranges 10--39\(^\circ\). This is structural: its sigmoid output scaling bounds
the magnitude head into a supplied band, and nothing anchors the angle. Its bus
schema also has no slot for bus voltage state, so it cannot observe an angle
datum.

\section{The pandapower GridSFM collapse was a data-conversion bug}
\label{sec:gridsfm}

Earlier revisions of this report treated GridSFM's frozen training as a property
of the model, and offered the leaky clip bounding its heads as the leading
explanation. \textbf{Both were wrong.} The cause has since been traced to the
parquet conversion, and correcting it makes GridSFM train on every pandapower
grid. The leaky-clip hypothesis is withdrawn.

\paragraph{Root cause: \texttt{S\_demand} is not the load.} In the ACOPF-backbone
parquets \texttt{S\_demand} (loaded as \texttt{S\_start}) stores the
\emph{pre-OPF injection} --- base generation minus scenario load --- not the
load. This was verified on case118 at pure-generator buses, where it equals the
base setpoint exactly: bus 9 at 450.00~MW, bus 24 at 220.00~MW, bus 25 at
314.00~MW.

The consequence is that \(S_{\rm opf} - S_{\rm demand}\) is a \textbf{redispatch
delta}, not a generation. Building the generation envelope on that delta puts
its range under the system-wide slack correction, which on case118 is about
\(-2271\)~MW. The resulting generation box has a \textbf{midpoint at
\(-28\)~pu}, which is not a generation level any model can use.

GridSFM builds its DC angle prior from the box midpoint. Fed a midpoint of
\(-28\)~pu it saturates immediately, and the gradient norm is then exactly zero
at every step --- which is precisely what the single-batch probe measured
(15{,}148{,}227 parameters with \texttt{requires\_grad=True}, a finite loss, a
connected autograd graph, and zero parameters receiving nonzero gradient). The
probe was correct; its interpretation was not. GridFM and LUMINA do not consume
the box midpoint, which is why the same malformed input degraded them without
freezing them.

\paragraph{The fix.} \texttt{convert\_opf\_backbone\_parquet.py -{}-gen\_box
injection} builds the envelope on \(S_{\rm opf}\) itself, which is observable and
physically bounded, while controllability is still decided by the delta --- a bus
is controllable iff the OPF actually moves its injection. Testing
\(|S_{\rm opf}|\) instead would flag every load bus and turn its power-balance
equality into a free variable. The flag reaches the models as the per-bus column
\texttt{Gen\_box\_injection}, read by \texttt{opf\_task.opf\_decision\_space}.

\paragraph{The confirming experiment.} Re-running the pandapower suite on the
corrected files reverses the failure completely:

\begin{table}[H]
\centering
\small
\begin{tabular}{lllll}
\toprule
Grid & Run & \multicolumn{2}{c}{original conversion} & corrected (\texttt{injbox}) \\
\cmidrule(lr){3-4}\cmidrule(lr){5-5}
 & & frozen? & train loss & train loss \\
\midrule
case14  & pre-FT  & no  & 1.0825e-02 (ep 40) & 2.3993e-03 (ep 35), moving \\
        & scratch & no  & 1.7211e-03 (ep 40) & 3.3299e-03 (ep 35), moving \\
case118 & pre-FT  & \textbf{yes} (ep 26--28 bit-identical) & 4.6775e-01 & 5.1315e-03 (ep 33), moving \\
        & scratch & \textbf{yes} (ep 23--25 bit-identical) & 3.9576e+00 & 1.3424e-01 (ep 30), moving \\
case300 & pre-FT  & no, but not converging & 5.6870e+00 & 1.4837e-01 (ep 29), moving \\
        & scratch & \textbf{yes} (ep 37--38) & 2.9180e+00 & 1.0378e-01 (ep 30), moving \\
\bottomrule
\end{tabular}
\caption{GridSFM on the original versus the corrected conversion. \textbf{The
decisive column is the third one.} Under the original conversion three runs are
frozen --- the training loss is bit-identical across consecutive epochs, which
is impossible under shuffled batches unless the weights are static. Under the
corrected conversion no bit-identical epoch appears on any grid: the freeze is
gone. The corrected runs were at epochs 29--35 of 40 at the time of writing, so
their final test numbers are not yet available.}
\end{table}

\paragraph{Loss magnitudes are not comparable across the two conversions.} The
loss contains a limit-violation term evaluated against the generation box, and
the box is exactly what changed, so the two loss columns above measure different
quantities. The large drops on case118 and case300 are mostly the disappearance
of an enormous violation penalty against a box centred at \(-28\)~pu, not
evidence that the model fits better. The clearest illustration is case14
scratch, where the corrected loss is \emph{higher} (1.72e-3 \(\rightarrow\)
3.33e-3): replacing a degenerate box with a real one gives the limit term
something genuine to penalise. Only the freeze/no-freeze column supports a
conclusion, and only the completed test metrics will support a capability claim.

\paragraph{What this invalidates.} The consequences run wider than the GridSFM
rows:

\begin{itemize}\itemsep2pt
\item \textbf{Every GridSFM row in Table~4 is an artifact of the bug}, not a
measurement of GridSFM. This is a stronger statement than earlier revisions
made: the rows are not merely ``runs that did not converge'', they are runs fed
a decision space with a generation box centred 28~pu away from any feasible
generation level.
\item \textbf{The pandapower cost gaps inherit the same defect}, since the cost
is evaluated against that box. This, and not an intrinsic instability of the
metric, is why the original and corrected cost gaps in Table~3 differ so wildly.
\item \textbf{GridFM and LUMINA's voltage columns largely survive}, because the
malformed box entered their loss only through the limit-violation term and never
through their architecture. Their corrected voltage RMSE moves by at most 13\%
(Table~3), which is why the model ranking is unchanged.
\end{itemize}

\paragraph{This is not a GridSFM pathology.} On the corrected pandapower data
GridSFM trains on all three grids; on OPFData, whose decision space is built from
the case file's true \texttt{pmin}/\texttt{pmax} and never passed through this
conversion, it trains on case14 and reaches 7.919e-3 pre-FT --- tied with LUMINA.
GridSFM's sensitivity here is a real and reportable property: it is the only one
of the three that consumes the generation box structurally, through its DC angle
prior, so a malformed box disables it outright where the others merely degrade.
That is a statement about coupling to input validity, not about capability.

\paragraph{A separate, milder phenomenon remains on OPFData.} Three of the four
OPFData GridSFM runs stop improving late in training: case14 scratch flattens at
9.7975e-03 around epoch 38, case118 scratch at 2.2444e-02 around epoch 31, and
case118 pre-FT diverges from 1.22e-01 to 2.6714e-01 and holds it from epoch 12.
Both case118 runs report a \(|V|\) RMSE of exactly 3.3899e-02 despite different
initialisations, and the pre-FT angle error sits at exactly \(20.000^\circ\),
which together indicate a magnitude head that has gone constant against a bound.

This is \emph{not} the mechanism described above --- OPFData's box is the true
generator box, correctly applied --- and it is much milder: it appears late,
after substantial convergence, rather than from the first step, and case14 pre-FT
is unaffected. It has not been diagnosed. Given that the analogous failure in the
other pipeline turned out to be a decision-space defect rather than a model
defect, the OPFData decision space is the first place to look, and no conclusion
about GridSFM should be drawn from these three runs in the meantime.

\section{Status and limitations}

\paragraph{Completeness.} A further 4 OPFData runs cover the released
\texttt{gridfm\_graphkit} model (2 grids \(\times\) 2 objectives,
Section~\ref{sec:graphkit}); all completed. All 20 earlier OPFData runs are
complete: 10 under the shared
objective and 10 under the native objectives, the latter run twice (alex and
helma). On the pandapower pipeline the corrected-conversion (\texttt{injbox})
suite is complete for GridFM and LUMINA on case14, case118 and case300, and for
GridSFM on case14 (both inits) and case118 pre-FT. Three GridSFM runs ---
case118 scratch and both case300 runs --- exceeded the 24\,h wall at epoch
37--38 of 40 and never reached their own test pass; they were scored afterwards
from their saved checkpoints and are marked \(\S\) in Table~5. SimBench has not
been re-converted, so its rows remain original-conversion only and are not
usable for GridSFM. Nothing in this report is now pending.

\paragraph{Two defects found while collecting these results.} Both are fixed and
both mattered. (i) The job label omitted the case, so case14 and case118 wrote
to the same SLURM stdout/stderr files and overwrote each other's diagnostics;
per-run training logs were unaffected. (ii) \texttt{-{}-EPOCHS 0} returns before
the post-training checkpoint load, so an attempt to score a saved checkpoint
silently evaluated the freshly initialised model instead --- it reported
plausible-looking numbers that were simply the untrained network. A
\texttt{-{}-resume\_state\_dict} flag was added to both drivers; it is the only
correct way to score a checkpoint this project produces, since
\texttt{-{}-init\_mode pretrained} routes through each repo's release loader and
rejects the plain \texttt{state\_dict} these scripts save.

\paragraph{Provenance of the cost column.} An earlier revision of this report
carried cost gaps for the pandapower case14 and SimBench runs. Those figures
could not be reproduced from the run logs --- those runs predate the masked OPF
metric block --- so they have been removed rather than restated. The pandapower
cost gaps that remain are quoted directly from the run logs, but they are
evaluated against the malformed generation box, and
Section~\ref{sec:metrics} explains why none of them should carry weight.

\paragraph{Blocked.} The LVN~Heo1 OPF dataset remains blocked:
\texttt{pp.runopp} does not converge on that CGMES-derived net even with the
voltage band opened to \([0,2]\) and all thermal limits removed, and the failure
survives switch fusion (3782\(\rightarrow\)1182 buses), zero-impedance repair,
and both initialisations. Progressing it needs a more robust backend
(PowerModels/IPOPT, which requires Julia) or repair of the CGMES import.

\paragraph{Not measured.} Branch ratings are stored in both pipelines but not fed
to the models; thermal violation is consequently not among the scored metrics and
all three models lack the feature equally. OPFData was run without topological
perturbations, so N-1 robustness is untested even though the dataset supports it.

\paragraph{The main open item.} Every number here comes from voltage-only
supervision. Both pipelines also carry the optimal dispatch and the objective,
neither of which was used as a target, and two of the three models (LUMINA,
GridSFM) expose generator heads that consequently received no gradient. A
dispatch-supervised variant is implemented (\texttt{opf\_task.dispatch\_loss}
and \texttt{gen\_head\_loss}, exposed as \texttt{-{}-dispatch\_weight} and
\texttt{-{}-gen\_head\_weight} on all training scripts) and should be run before
these results are used to rank the models on anything but voltage accuracy. It
was investigated as a possible fix for GridSFM's frozen training and is not one
--- a single-batch test shows gradients still do not reach \texttt{head\_Pg} when
it is supervised directly. That question is now moot for the pandapower pipeline,
where the freeze had a data cause and is resolved (Section~\ref{sec:gridsfm}).

\end{document}
